\sscp{Austronesian incursion/migration(s): \tsc{Aru}}{15-06-08}
It is not yet clear where Austronesian-speaking peoples came
from to start what eventually developed into the \tsc{Aru} subgroup (see \crf{ch:Aru}).
These languages show strong substrate influence.
Conservative features such as the retention of PMP *tuqəlaŋ ‘bone’
(rather than reflexes of *duRi ‘bone’) certainly suggest a primary,
rather than secondary migration.
But everything is masked by heavy substrate influence
and a large replacement of basic vocabulary.
Where did the first group of Austronesian speakers come from?
We note in many distribution maps throughout this study that
\tsc{Aru} is often different in noticeable ways.

The geographically obvious possibility is via the neighbouring
\tsc{Seram-Tanimbar-Bomberai} subgroup, most members of which share the common sound change of *p > **f/ɸ with \tsc{Aru}.
But it is difficult to reconcile other sound correspondences
to clearly establish such a connection with available data.
\tsc{Timor-Babar} would be another possibility,
which also probably shares the very common sound change of *p
> **f with \tsc{Aru}.
But that has the same problem of being difficult to reconcile
the other sound correspondences to clearly establish such a connection with available data.
The possibility of linking the \tsc{Aru} languages within “Central
Maluku” suggested \citet{GrayGreenhill2009} appears to have little to support it in the historical
phonology or the lexicon.

The idea of a link further east deserves further investigation,
as shown by the fact that during the many years that Grimes was
resident in Maluku, there was a case of a small dingy with two men
and a woman who had run out of fuel in the Torres Straits,
and then drifted on the currents to Aru with one surviving to
reach those islands (pers. comm.
Jock Hughes, confirmed by email 30 October 2020).
So a possible (secondary) link to the east with Austronesian
languages such as Motu or Papuan Tip languages should not be
ruled out without investigation. See also \citet{Voorhoeve1982}.